\documentclass[10pt]{article}
\usepackage[T1]{fontenc}
\usepackage[utf8]{inputenc}
\usepackage[french]{babel}
\usepackage{amsmath,amssymb,mathrsfs,mathpazo}
\usepackage{geometry}
\geometry{a4paper,margin=2cm}
\usepackage{graphicx,array,tabularx,tabularray,colortbl,multirow,multicol,booktabs,enumitem,pifont,ulem,fancybox,diagbox,cancel}
\usepackage{pstricks,pstricks-add,pst-plot,pst-tree,pst-node,pst-text,pst-func,pst-math,pst-bspline,pst-3dplot,pst-eucl,pst-solides3d}
\let\pscalculate\undefined
\usepackage{pst-bezier}
\usepackage{tikz,pgfplots,tkz-tab}
\usetikzlibrary{arrows,arrows.meta,patterns,calc,positioning,trees,decorations.markings,shapes.geometric,shapes.misc,shapes.arrows,intersections,decorations.pathmorphing,decorations.pathreplacing,backgrounds,shadows,babel}
\pgfplotsset{compat=1.18}
\usepgfplotslibrary{fillbetween}
\usepackage[np]{numprint}
\usepackage{listings,makecell,pas-tableur,twoopt}
\let\barre\undefined
\usepackage[pstricks]{tablvar}
\definecolor{darkgreen}{RGB}{0,100,0}
\providecommand{\red}{\color{red}}
\providecommand{\blue}{\color{blue}}
\providecommand{\green}{\color{green}}
\providecommand{\black}{\color{black}}
\providecommand{\white}{\color{white}}
\providecommand{\gray}{\color{gray}}
\providecommand{\cyan}{\color{cyan}}
\providecommand{\magenta}{\color{magenta}}
\providecommand{\yellow}{\color{yellow}}
\providecommand{\violet}{\color{violet}}
\providecommand{\orange}{\color{orange}}
\providecommand{\euro}{\texteuro}
\providecommand{\EUR}{\texteuro}
\providecommand{\ieme}{\textsuperscript{e}}
\providecommand{\degres}{\ensuremath{^\circ}}
\providecommand{\degree}{\ensuremath{^\circ}}
\providecommand{\vv}[1]{\overrightarrow{#1}}
\providecommand{\Coord}[1]{#1}
\providecommand{\decosix}{}
\providecommand{\starredbullet}{\ensuremath{\bullet}}
\pagestyle{empty}
\setlength{\parindent}{0pt}
\usepackage{adjustbox}\DeclareUnicodeCharacter{2236}{:}\DeclareUnicodeCharacter{2190}{\ensuremath{\leftarrow}}
\providecommand{\R}{\mathbb{R}}
\providecommand{\N}{\mathbb{N}}
\providecommand{\Z}{\mathbb{Z}}
\providecommand{\Q}{\mathbb{Q}}
\providecommand{\D}{\mathbb{D}}
\providecommand{\C}{\mathbb{C}}
\providecommand{\e}{{\mathrm{e}}}
\providecommand{\vect}[1]{\overrightarrow{#1}}
\providecommand{\vectt}[1]{\overrightarrow{\mathrm{#1}}}
\providecommand{\vvt}[1]{\overrightarrow{\mathrm{#1}}}
\providecommand{\np}[2][]{{#2~#1}}
\providecommand{\vv}[1]{{\overrightarrow{#1}}}
\renewcommand{\d}{{\mathrm{d}}}
\providecommand{\degres}{{^\circ}}
\providecommand{\degre}{{^\circ}}
\providecommand{\blue}{\color{blue}}
\providecommand{\red}{\color{red}}
\providecommand{\enskip}{\quad}
\providecommand{\llbracket}{\lbrack\!\lbrack}
\providecommand{\rrbracket}{\rbrack\!\rbrack}
\providecommand{\euro}{{\text{€}}}
\providecommand{\no}{{\text{n°}}}
\providecommand{\interval}[2]{{[#1;#2]}}
\providecommand{\intervFO}[2]{{[#1;#2[}}
\providecommand{\dotfill}{\ldots}

\providecommand{\ds}{\displaystyle}
\providecommand{\textding}[1]{#1}
\providecommand{\Oij}{\(\mathrm{O};\vec{\imath},\vec{\jmath}\)}
\providecommand{\Oijk}{\(\mathrm{O};\vec{\imath},\vec{\jmath},\vec{k}\)}
\providecommand{\Ouv}{\(\mathrm{O};\vec{u},\vec{v}\)}
\providecommand{\intervFF}[2]{\ensuremath{[#1;#2]}}
\renewcommand{\euro}{\ifmmode\text{\texteuro}\else\texteuro\fi}
\DeclareUnicodeCharacter{00B0}{\ifmmode^{\circ}\else\textdegree\fi}
\DeclareUnicodeCharacter{20AC}{\ifmmode\text{\texteuro}\else\texteuro\fi}
% Énoncés : Éducation nationale. Transcription et corrigés : APMEP, Denis Vergès et contributeurs. Découpage et classement : CoopMaths, Sébastien Lozano, Éric Elter et contributeurs. Adaptation de la présentation : Maths Explorer.
% CC BY-SA : https://coopmaths.fr/www/about/
% https://coopmaths.fr/alea/static/bac/2023/tex/bac_2023_03_sujet2_etrangers_groupe1_1.tex
\begin{document}
\noindent\textit{APMEP / CoopMaths. CC BY-SA. Présentation adaptée par Maths Explorer.}\par\medskip
\medskip

Cet exercice est un questionnaire à choix multiples. 

Pour chaque question, une seule des quatre propositions est exacte. 

Indiquer sur la copie le numéro de la question et la lettre de la proposition choisie. 

Aucune justification n'est demandée.

Pour chaque question, une réponse exacte rapporte un point. Une réponse fausse, une réponse multiple ou l'absence de réponse ne rapporte ni n'enlève de point.

\medskip

\textbf{Question 1 :}

\medskip

Soit $f$ la fonction définie sur $\R$ par $f(x) = x\text{e}^x$.

Une primitive $F$ sur $\R$ de la fonction $f$ est définie par:
\begin{center}
\begin{tabularx}{\linewidth}{X X}
\textbf{A.} $F(x) = \dfrac{x^2}{2}\text{e}^x$&\textbf{B} $F(x) = (x - 1)\text{e}^x$\\
\textbf{C.} $F(x) = (x + 1)\text{e}^x$&\textbf{D.} $F(x) = \dfrac2 x \text{e}^{x^2}$.
\end{tabularx}
\end{center}

\medskip

\textbf{Question 2 :}

\medskip

La courbe $\mathcal{C}$ ci-dessous représente une fonction $f$ définie et deux fois dérivable
sur $]0~;~+\infty[$. On sait que :

\begin{itemize}
\item[$\bullet~~$] le maximum de la fonction $f$ est atteint au point d'abscisse 3 ;
\item[$\bullet~~$] le point P d'abscisse 5 est l'unique point d'inflexion de la courbe $\mathcal{C}$.
\end{itemize}

\begin{center}
\psset{unit=0.85cm,arrowsize=2pt 4}
\begin{pspicture*}(-1,-3)(14,6)
\psaxes[linewidth=1.25pt,labelFontSize=\scriptstyle]{->}(0,0)(-1,-3)(14,6)
\psplot[plotpoints=2000,linewidth=1.25pt,linecolor=red]{0}{14}{x 1 sub 2.71828 x 0.5 mul exp div 10 mul}
\psplotTangent[linestyle=dashed]{5}{5}{x 1 sub 2.71828 x 0.5 mul exp div 10 mul}
\psplotTangent[linestyle=dashed]{3}{2}{x 1 sub 2.71828 x 0.5 mul exp div 10 mul}
\uput[r](0.85,-2){\red $\mathcal{C}$}
\psdots[dotstyle=+,dotangle=45,dotscale=1.5](5,3.2834)\uput[ur](5,3.2834){P}
\end{pspicture*}
\end{center}

On a:

\begin{center}
\begin{tabularx}{\linewidth}{X X}
\textbf{A.~~} pour tout $x \in  ]0~;~5[,\:f(x)$ et $f'(x)$ sont de même signe;&
\textbf{B.~~} pour tout $x \in ]5~;~+\infty[\:, f(x)$ et $f'(x)$ sont de même signe;\\
\textbf{C.~~} pour tout $x \in ]0~;~5[,\: f'(x)$ et $f''(x)$ sont de même signe;&
\textbf{D.~~} pour tout $x \in  ]5~;~+\infty[,\: f(x)$ et $f''(x)$ sont de même signe.
\end{tabularx}
\end{center}

\medskip

\textbf{Question 3 :}

\medskip

On considère la fonction $g$ définie sur $[0~;~ +\infty[$ par $g(t) = \dfrac{a}{b + \text{e}^{-t}}$ où $a$ et $b$ sont deux nombres réels.

On sait que $g(0) = 2$ et $\displaystyle\lim_{t \to + \infty}  g(t) = 3$. 

Les valeurs de $a$ et $b$ sont:

\begin{center}
\begin{tabularx}{\linewidth}{X X}
\textbf{A.~~} $a = 2$ et $b = 3$&\textbf{B.~~} $a = 4$ et $b = \dfrac43$\\
\textbf{C.~~} $a = 4$ et $b = 1$&\textbf{D.~~}$a = 6$ et $b = 2$
\end{tabularx}
\end{center}

\medskip

\textbf{Question 4 :}

\medskip

Alice dispose de deux urnes A et B contenant chacune quatre boules indiscernables au toucher.

L'urne A contient deux boules vertes et deux boules rouges.

L'urne B contient trois boules vertes et une boule rouge.

Alice choisit au hasard une urne puis une boule dans cette urne. Elle obtient une boule verte.

La probabilité qu'elle ait choisi l'urne B est:

\begin{center}
\renewcommand\arraystretch{2.1}
\begin{tabularx}{\linewidth}{X X}
\textbf{A.~~} $\dfrac38$&\textbf{B.~~} $\dfrac12$\\
\textbf{C.~~} $\dfrac35$&\textbf{D.~~} $\dfrac58$
\end{tabularx}
\end{center}
\renewcommand\arraystretch{1}

\medskip

\textbf{Question 5 :}

\medskip

On pose $S = 1 + \dfrac12  + \dfrac13 + \dfrac14 + \ldots + \dfrac{1}{100}$.

Parmi les scripts Python ci-dessous, celui qui permet de calculer la somme $S$ est :

\begin{center}
\begin{tabularx}{\linewidth}{X X}
\textbf{a.~~}&\textbf{b.~~}\\
\begin{tabular}{l}
def somme$\_{}$a():\\
\quad S = 0\\
\quad for k in range(100):\\
\qquad S=1/(k+1) \\
\quad return S
\end{tabular}&\begin{tabular}{l}
def somme$\_{}$b():\\
\quad S = 0\\
\quad for k in range(100) :\\
\qquad S = S + 1/(k + 1)\\
\quad return S
\end{tabular}\\
\textbf{c.~~}&\textbf{d.~~}\\
\begin{tabular}{l}
def somme$\_{}$c():\\
\quad k = 0\\
\quad while S < 100 :\\
\qquad S = S+1/(k+1)\\
\quad return S
\end{tabular}&
\begin{tabular}{l}
def somme$\_{}$d():\\
\quad k = 0\\
\quad while k $< 100$: \\
\qquad S = S + 1/(k + 1)\\
\quad return S
\end{tabular}
\end{tabularx}
\end{center}

\bigskip

\end{document}
